%% kapa: exam=1 round=36 year=2025 period=1 form=A subject=economics subject_ko=경제학원론 seq=21 number=61
%% source: Q-Net 감정평가사 기출문제 (cst00309) artlSeq=5231525 file=raw/1008/5231525/2025년 제36회 감정평가사 1차시험 1교시 문제지 원본.pdf
%% conversion: kordoc 4.24.51 -> tools/gen1.py (auto); review=proofread-sonnet; flags=none
\begin{question}{61}
유동성선호설을 통한 화폐시장과 LM곡선에 대해 옳게 설명한 것은? (단, 화폐공급곡선은 수직이고, 화폐수요는 이자율에 음(-), 소득에 정(+)의 관계를 가정)
\begin{choices}
  \item 소득의 감소는 화폐수요곡선을 오른쪽으로 이동시켜 이자율 상승
  \item 소득의 증가는 실질화폐 수요를 증가시켜 이자율 상승
  \item 이자율이 상승하면 화폐수요곡선이 왼쪽으로 이동하고 소득 증가
  \item 중앙은행이 화폐공급을 감소하면 이자율 하락
  \item 투기적 동기로 인한 소득의 증가는 화폐수요곡선상에서의 이동으로 이자율 상승
\end{choices}
\end{question}
\answer{2}
